\documentclass[11pt]{article}
\usepackage[T1]{fontenc}
\usepackage{lmodern}
\usepackage{microtype}
\usepackage[margin=1in]{geometry}
\usepackage{amsmath,amssymb,amsthm,mathtools}
\usepackage{mathrsfs}
\usepackage{booktabs,array,longtable}
\usepackage{enumitem}
\usepackage{tikz-cd}
\usepackage[hidelinks]{hyperref}
\usepackage[nameinlink,noabbrev]{cleveref}
\usepackage{xcolor}

\setlist[itemize]{leftmargin=2em,itemsep=2pt,topsep=4pt}
\setlist[enumerate]{leftmargin=2.2em,itemsep=2pt,topsep=4pt}
\allowdisplaybreaks

\newtheorem{theorem}{Theorem}[section]
\newtheorem{proposition}[theorem]{Proposition}
\newtheorem{lemma}[theorem]{Lemma}
\newtheorem{corollary}[theorem]{Corollary}
\theoremstyle{definition}
\newtheorem{definition}[theorem]{Definition}
\newtheorem{axiom}[theorem]{Axiom}
\newtheorem{construction}[theorem]{Construction}
\newtheorem{example}[theorem]{Example}
\newtheorem{conjecture}[theorem]{Conjecture}
\newtheorem{problem}[theorem]{Open Problem}
\theoremstyle{remark}
\newtheorem{remark}[theorem]{Remark}

\newcommand{\K}{\mathbf K}
\newcommand{\Rcat}{\mathbf R}
\newcommand{\Tcat}{\mathbf T}
\newcommand{\Cat}{\mathbf{Cat}}
\newcommand{\Ob}{\operatorname{Ob}}
\newcommand{\Hom}{\operatorname{Hom}}
\newcommand{\Lift}{\operatorname{Lift}}
\newcommand{\Cofib}{\operatorname{Cofib}}
\newcommand{\Recon}{\operatorname{Recon}}
\newcommand{\Obs}{\operatorname{Obs}}
\newcommand{\Zoo}{\mathcal Z}
\newcommand{\Seeds}{\Sigma}
\newcommand{\Res}{\mathsf{Res}}
\newcommand{\Trans}{\mathsf{Trans}}
\newcommand{\id}{\operatorname{id}}
\newcommand{\colim}{\operatorname*{colim}}
\newcommand{\status}[1]{\textbf{Status: #1.}}

\title{\textbf{Komorebi Seeded Two-Axis Calculus}\\
\large Transformation, Resolution, Residual Generation, and Reconstruction\\[0.4em]
\normalsize Working Formalization v0.5.0}
\author{Komorebi P.}
\date{26 September 2026}

\begin{document}
\maketitle

\begin{abstract}
This note develops a formal version of the Komorebi seed calculus in which two kinds of mathematical motion are treated as distinct primitives: transformation and resolution. Transformation changes representation inside a fixed level of distinguishability. Resolution changes which distinctions are visible. The basic architecture is an indexed category $F:\Rcat^{op}\to\Cat$ and its Grothendieck construction $\pi:\K\to\Rcat$. This already yields a canonical factorization of mixed morphisms into a within-resolution transformation followed by a cartesian resolution lift. Seeds are then closed under transformation, coarsening, admissible refinement, and residual generation. A monotone closure operator gives a least generated ``zoo.'' Exchange maps measure whether transformation commutes with resolution. In stable or abelian fibers their cofibers, kernels, or cokernels define mixed residuals. Several basic results are proved here: resolution monotonicity, existence and factorization of the total category, existence of least seed zoos, a resolution ultrametric under nested finite resolutions, lifting of resolution-compatible transformations to the total category, and an exact normal-form theorem for terminating confluent exchange systems. Version 0.5.0 closes two additional stable-layer obligations, adds a nonstable descent test bed, and proves an obstruction-localization principle. A normalization path in a stable enhancement has a canonical total defect filtration whose successive quotients are transported elementary exchange residuals. The obstruction layer is represented by a Verdier quotient of the stable seed zoo by the thick closure of operationally killable residuals; under a directed common-annihilator hypothesis, operational killability itself is thick and is detected exactly by the quotient. Exactness then forces the total normalization obstruction to vanish whenever all elementary exchange obstructions vanish, and any persistent total obstruction localizes to at least one elementary exchange. Separately, the local mismatch data of an abelian torsor form a Cech residual cocycle whose class modulo local repair is precisely the obstruction to reconstructing a global section. Later sections retain the stronger order-independent, arithmetic reconstruction, and local-global reciprocity problems. The note is deliberately explicit about which statements are elementary consequences of standard categorical machinery and which are proposed Komorebi-specific extensions.
\end{abstract}

\tableofcontents
\newpage

\section{Purpose and status}

The goal is to formalize the idea that a mathematical theory may contain two distinct ambient directions:
\[
\boxed{\text{transformation: change of representation}}
\qquad\text{and}\qquad
\boxed{\text{resolution: change of distinguishability}.}
\]
The distinction is operational rather than rhetorical. A Fourier transform, a duality, or a change of coordinates may alter representation while preserving the level of information being resolved. Coarsening a filtration, passing from modulus $p^{n+1}$ to modulus $p^n$, truncating a spectral expansion, or forgetting fine ramification data changes the level at which distinctions are visible.

The central hypothesis of this note is that a large class of mathematical constructions can be organized by a two-axis calculus in which:
\begin{enumerate}
    \item objects occur at resolutions;
    \item each resolution has its own transformation category;
    \item resolution changes induce functors between those categories;
    \item mixed transformation-resolution squares can commute exactly, commute up to equivalence, or leave a structured residual;
    \item residuals can themselves be promoted to new seeds;
    \item compatible local or resolved data can sometimes be reconstructed into global objects.
\end{enumerate}

\paragraph{Status convention.}
Results tagged ``proved here'' are elementary theorems under the stated hypotheses. Results tagged ``standard machinery'' are direct applications of established category, order, or rewriting theory. Statements tagged ``program theorem'' are precise targets whose Komorebi-specific hypotheses remain to be supplied in applications. No claim of mathematical novelty is made merely because standard machinery is expressed in Komorebi language.

\section{Resolution systems}

\subsection{The resolution base}

\begin{definition}[Resolution category]
A \emph{resolution category} is a small poset $(R,\le)$ regarded as a thin category $\Rcat$. An arrow $r\to s$ exists exactly when $r\le s$. We interpret $s$ as at least as fine as $r$.
\end{definition}

The poset assumption is intentionally restrictive for the first formalization. A later version may replace $\Rcat$ by a general category, a site, a filtered category, or a multi-parameter resolution geometry.

\begin{definition}[Indexed transformation system]
An \emph{indexed transformation system} over $\Rcat$ is a strict functor
\[
F:\Rcat^{op}\longrightarrow\Cat.
\]
For $r\in R$, write
\[
\Tcat_r:=F(r).
\]
For $r\le s$, write
\[
C_{s\to r}:=F(r\le s):\Tcat_s\longrightarrow\Tcat_r.
\]
We call $C_{s\to r}$ the \emph{coarsening functor}.
\end{definition}

Strictness gives
\[
C_{r\to r}=\id_{\Tcat_r}
\]
and, for $r\le s\le t$,
\[
C_{t\to r}=C_{s\to r}\circ C_{t\to s}.
\]
A pseudofunctor version is more natural in some applications. The strict version isolates the main structure without coherence bookkeeping.

\begin{definition}[Resolved object]
A \emph{resolved object} is a pair $(r,X)$ with $X\in\Ob(\Tcat_r)$.
\end{definition}

\subsection{Resolution equivalence}

Let $X,Y\in\Tcat_t$ and let $r\le t$.

\begin{definition}[Indistinguishability at resolution $r$]
We write
\[
X\sim_r Y
\]
when
\[
C_{t\to r}(X)\cong C_{t\to r}(Y)
\]
in $\Tcat_r$.
\end{definition}

For a fixed ambient fine level $t$, each $\sim_r$ is an equivalence relation on isomorphism classes of objects of $\Tcat_t$.

\begin{theorem}[Resolution monotonicity]\label{thm:resolution-monotonicity}
Let $r\le s\le t$ and $X,Y\in\Tcat_t$. If
\[
X\sim_sY,
\]
then
\[
X\sim_rY.
\]
\end{theorem}

\begin{proof}
An isomorphism
\[
C_{t\to s}(X)\cong C_{t\to s}(Y)
\]
is carried by the functor $C_{s\to r}$ to an isomorphism
\[
C_{s\to r}C_{t\to s}(X)
\cong
C_{s\to r}C_{t\to s}(Y).
\]
By functorial coherence this is
\[
C_{t\to r}(X)\cong C_{t\to r}(Y).
\]
Hence $X\sim_rY$.
\end{proof}

\begin{corollary}[Nested distinguishability]\label{cor:nested}
For $r\le s\le t$, the equivalence relation $\sim_s$ refines $\sim_r$.
\end{corollary}

This is the first exact form of the statement that coarsening may erase distinctions but cannot create distinctions that were absent at finer resolution.

\subsection{Refinement as a fiber rather than a function}

\begin{definition}[Lift set]
For $r\le s$ and $X\in\Tcat_r$, define the set of isomorphism classes of lifts
\[
\Lift_{r\to s}(X)
:=
\left\{
Y\in\Ob(\Tcat_s):C_{s\to r}(Y)\cong X
\right\}/\cong.
\]
\end{definition}

Three basic cases occur:
\[
|\Lift_{r\to s}(X)|=0,
\qquad
|\Lift_{r\to s}(X)|=1,
\qquad
|\Lift_{r\to s}(X)|>1.
\]
They correspond respectively to obstruction to refinement, unique refinement, and hidden fine-scale multiplicity.

\begin{definition}[Resolution branching number]
When $\Lift_{r\to s}(X)$ is finite, define
\[
b_{r\to s}(X):=|\Lift_{r\to s}(X)|.
\]
Its logarithm
\[
H_{r\to s}(X):=\log b_{r\to s}(X)
\]
may be regarded as a crude \emph{resolution information} measuring how much ambiguity is hidden by coarsening.
\end{definition}

This quantity is only a first model. In probabilistic, measurable, or derived settings one would replace cardinality by entropy, measure, dimension, rank, or a homotopical invariant.

\section{The total two-axis category}

\subsection{Grothendieck construction}

\begin{construction}[Total Komorebi category]\label{cons:groth}
Given $F:\Rcat^{op}\to\Cat$, define the category
\[
\K:=\int_{\Rcat}F
\]
as follows.

Its objects are pairs $(r,X)$ with $X\in\Tcat_r$.

For $r\le s$, a morphism
\[
(r,X)\longrightarrow(s,Y)
\]
is a morphism
\[
\alpha:X\longrightarrow C_{s\to r}(Y)
\]
in $\Tcat_r$. If $r\nleq s$, there is no such morphism in this directed version.
\end{construction}

Given
\[
(r,X)\xrightarrow{\alpha}(s,Y)
\quad\text{and}\quad
(s,Y)\xrightarrow{\beta}(t,Z),
\]
define
\[
\beta\circ_{\K}\alpha
:=
C_{s\to r}(\beta)\circ\alpha:
X\to C_{t\to r}(Z).
\]

\begin{theorem}[Total category]\label{thm:total-category}
The construction above defines a category $\K$.
\end{theorem}

\begin{proof}
Identity at $(r,X)$ is $\id_X:X\to C_{r\to r}(X)=X$. For associativity, let
\[
(r,X)\xrightarrow{\alpha}(s,Y)
\xrightarrow{\beta}(t,Z)
\xrightarrow{\gamma}(u,W).
\]
One composite gives
\[
C_{t\to r}(\gamma)\circ C_{s\to r}(\beta)\circ\alpha.
\]
The other gives the same expression by functoriality and
\[
C_{t\to r}=C_{s\to r}C_{t\to s}.
\]
Thus composition is associative.
\end{proof}

Define
\[
\pi:\K\to\Rcat,
\qquad
\pi(r,X)=r.
\]
The fiber over $r$ is canonically equivalent to $\Tcat_r$.

\subsection{Canonical factorization}

For $r\le s$ and $Y\in\Tcat_s$, define the canonical morphism
\[
\kappa_{r,s,Y}:
(r,C_{s\to r}Y)\longrightarrow(s,Y)
\]
represented by the identity map
\[
\id_{C_{s\to r}Y}.
\]

\begin{theorem}[Two-axis factorization]\label{thm:factorization}
Every morphism
\[
f:(r,X)\longrightarrow(s,Y)
\]
in $\K$ factors uniquely as
\[
(r,X)
\xrightarrow{f_v}
(r,C_{s\to r}Y)
\xrightarrow{\kappa_{r,s,Y}}
(s,Y),
\]
where $f_v$ lies entirely in the fiber $\Tcat_r$.
\end{theorem}

\begin{proof}
By definition, $f$ is exactly a morphism
\[
\alpha:X\to C_{s\to r}Y
\]
in $\Tcat_r$. Take $f_v=\alpha$. Composition with $\kappa_{r,s,Y}$ is represented by
\[
\id_{C_{s\to r}Y}\circ\alpha=\alpha,
\]
so the composite is $f$. Any factorization through the fixed cartesian target $(r,C_{s\to r}Y)$ must use the same fiber morphism $\alpha$, hence is unique.
\end{proof}

\begin{remark}
This is the formal version of the two-axis statement
\[
\boxed{
\text{mixed movement}
=
\text{within-resolution transformation}
+
\text{resolution movement}.
}
\]
Categorically, $\pi:\K\to\Rcat$ is the Grothendieck fibration associated with $F$.
\end{remark}

\section{Resolution geometry}

The nested equivalence relations from \cref{cor:nested} induce a genuine geometry when the resolution set is linearly ordered.

\begin{definition}[Separation depth]
Assume $R=\mathbb N$ with larger $n$ finer. Let $X,Y$ be objects represented at all sufficiently fine levels and let $\sim_n$ be nested equivalence relations. Define
\[
\nu(X,Y):=\sup\{n\in\mathbb N:X\sim_nY\}\in\mathbb N\cup\{\infty\}.
\]
Assume the system separates points, meaning $\nu(X,Y)=\infty$ only when $X=Y$ in the chosen underlying set of states.
\end{definition}

\begin{theorem}[Resolution ultrametric]\label{thm:ultrametric}
Define
\[
d_R(X,Y):=
\begin{cases}
0,&X=Y,\\
2^{-\nu(X,Y)},&X\ne Y.
\end{cases}
\]
Then $d_R$ is an ultrametric:
\[
d_R(X,Z)\le\max\{d_R(X,Y),d_R(Y,Z)\}.
\]
\end{theorem}

\begin{proof}
Let
\[
m=\min\{\nu(X,Y),\nu(Y,Z)\}.
\]
For every $n\le m$, we have $X\sim_nY$ and $Y\sim_nZ$. Since $\sim_n$ is an equivalence relation, $X\sim_nZ$. Hence
\[
\nu(X,Z)\ge m.
\]
Therefore
\[
2^{-\nu(X,Z)}
\le
2^{-m}
=
\max\{2^{-\nu(X,Y)},2^{-\nu(Y,Z)}\}.
\]
The identity and symmetry axioms are immediate.
\end{proof}

Thus the resolution axis can carry an intrinsic non-Archimedean geometry. Two objects are close when many levels of resolution are required before they can be distinguished.

\begin{definition}[First distinguishing scale]
When finite, define
\[
\delta(X,Y):=\nu(X,Y)+1.
\]
This is the first resolution at which $X$ and $Y$ become distinguishable.
\end{definition}

This quantity may be useful in $p$-adic, dyadic, symbolic, filtration, and machine-representation settings.

\section{Seeds and generated zoos}

\subsection{Seed signatures}

Let $U$ be a set of isomorphism classes of resolved objects under consideration. Smallness is imposed only to use ordinary powerset fixed-point arguments.

\begin{definition}[Seed signature]
A \emph{seed signature} on $U$ is data
\[
\mathfrak S=(\Seeds_0,\mathscr T,\mathscr C,\mathscr L,\mathscr D),
\]
where:
\begin{itemize}
    \item $\Seeds_0\subseteq U$ is the initial seed set;
    \item $\mathscr T$ is a collection of partial unary transformation constructors on $U$;
    \item $\mathscr C$ is the collection of allowed coarsening constructors;
    \item $\mathscr L$ assigns to an object its admissible refinement classes;
    \item $\mathscr D$ is a collection of partial residual constructors.
\end{itemize}
\end{definition}

The use of partial constructors is intentional. An operation need not be defined on every object or every resolution.

\begin{definition}[One-step zoo operator]
For $A\subseteq U$, define
\[
\Phi(A):=
A
\cup \mathscr T[A]
\cup \mathscr C[A]
\cup \mathscr L[A]
\cup \mathscr D[A],
\]
where each term denotes all defined outputs of the corresponding constructors applied to inputs from $A$.
\end{definition}

\begin{lemma}[Monotonicity]\label{lem:phi-monotone}
If $A\subseteq B$, then
\[
\Phi(A)\subseteq\Phi(B).
\]
\end{lemma}

\begin{proof}
Every constructor output available from inputs in $A$ is also available from the same inputs viewed as elements of $B$.
\end{proof}

\begin{theorem}[Seed zoo existence]\label{thm:zoo-existence}
There exists a least $\Phi$-closed set containing $\Seeds_0$. Equivalently, there is a least fixed point of the closure process above $\Seeds_0$.
\end{theorem}

\begin{proof}
The powerset $\mathcal P(U)$ is a complete lattice. The map
\[
\Psi(A):=\Seeds_0\cup\Phi(A)
\]
is monotone by \cref{lem:phi-monotone}. By the Knaster--Tarski fixed-point theorem, $\Psi$ has a least fixed point. This is exactly the least constructor-closed set containing $\Seeds_0$.
\end{proof}

\begin{definition}[Generated zoo]
The least fixed point in \cref{thm:zoo-existence} is denoted
\[
\Zoo(\Seeds_0).
\]
\end{definition}

\begin{proposition}[Finitary iteration]\label{prop:finitary}
Assume every constructor uses only finitely many inputs and $\Phi$ preserves increasing unions of $\omega$-chains. Then
\[
\Zoo(\Seeds_0)
=
\bigcup_{n<\omega}\Psi^n(\varnothing).
\]
\end{proposition}

\begin{proof}
The union on the right is a fixed point because $\Psi$ preserves the increasing union. It contains $\Seeds_0$ and is contained in every fixed point by induction on $n$, hence is least.
\end{proof}

This turns the seed calculus into an explicit generative process:
\[
\Seeds_0
\subseteq
\Seeds_1
\subseteq
\Seeds_2
\subseteq\cdots,
\qquad
\Seeds_{n+1}=\Psi(\Seeds_n).
\]

\subsection{Zoo growth}

\begin{definition}[Growth function]
When each finite stage contains finitely many isomorphism classes, define
\[
g_{\Seeds_0}(n):=|\Psi^n(\varnothing)|.
\]
\end{definition}

The asymptotic class of $g(n)$ may distinguish finite, polynomial, exponential, and faster-growing seed theories. This is a proposed invariant of a generated mathematical grammar rather than a theorem of the present note.

\section{Transformation-resolution exchange}

\subsection{Resolution-indexed transformations}

\begin{definition}[Lax resolution-compatible transformation]\label{def:laxtrans}
A \emph{lax resolution-compatible transformation} consists of endofunctors
\[
A_r:\Tcat_r\to\Tcat_r
\]
for every $r\in R$, together with natural transformations
\[
\eta^A_{s,r}:
A_rC_{s\to r}
\Longrightarrow
C_{s\to r}A_s
\qquad(r\le s),
\]
satisfying identity and composition coherence:
\[
\eta^A_{r,r}=\id
\]
and, for $r\le s\le t$,
\[
\eta^A_{t,r}
=
\bigl(C_{s\to r}\eta^A_{t,s}\bigr)
\circ
\bigl(\eta^A_{s,r}C_{t\to s}\bigr).
\]
If every $\eta^A_{s,r}$ is a natural isomorphism, $A$ is called \emph{exactly resolution-compatible}.
\end{definition}

The square is
\[
\begin{tikzcd}[column sep=large,row sep=large]
A_rC_{s\to r}(X) \arrow[r,"\eta^A_{s,r,X}"] & C_{s\to r}A_s(X).
\end{tikzcd}
\]
It compares transforming after coarsening with coarsening after transforming.

\begin{theorem}[Lift to the total category]\label{thm:lift-total}
Every lax resolution-compatible transformation $A$ induces an endofunctor
\[
\widehat A:\K\to\K
\]
defined on objects by
\[
\widehat A(r,X)=(r,A_rX).
\]
\end{theorem}

\begin{proof}
A morphism $(r,X)\to(s,Y)$ in $\K$ is represented by
\[
\alpha:X\to C_{s\to r}Y.
\]
Define $\widehat A(\alpha)$ to be
\[
A_rX
\xrightarrow{A_r(\alpha)}
A_rC_{s\to r}Y
\xrightarrow{\eta^A_{s,r,Y}}
C_{s\to r}A_sY.
\]
Identity preservation follows from $\eta^A_{r,r}=\id$. Composition preservation is exactly the coherence condition in \cref{def:laxtrans}, together with naturality of $\eta^A$.
\end{proof}

Thus a family of transformations defined independently at each resolution becomes a single transformation of the whole two-axis system precisely when its exchange data are coherent.

\subsection{Mixed residuals}

\begin{definition}[Stable mixed residual]
Assume each $\Tcat_r$ is stable and each coarsening functor is exact. For $r\le s$ and $X\in\Tcat_s$, define
\[
\Delta^A_{s,r}(X)
:=
\Cofib\bigl(\eta^A_{s,r,X}\bigr)
\in\Tcat_r.
\]
\end{definition}

Then there is a cofiber sequence
\[
A_rC_{s\to r}X
\longrightarrow
C_{s\to r}A_sX
\longrightarrow
\Delta^A_{s,r}(X).
\]

\begin{proposition}[Residual vanishing criterion]\label{prop:residual-vanish}
In a stable fiber,
\[
\Delta^A_{s,r}(X)\simeq0
\]
if and only if
\[
\eta^A_{s,r,X}
\]
is an equivalence.
\end{proposition}

\begin{proof}
In a stable category a morphism is an equivalence exactly when its cofiber is a zero object.
\end{proof}

In an abelian fiber one may instead record
\[
\Delta^-_{s,r}(X)=\ker\eta^A_{s,r,X},
\qquad
\Delta^+_{s,r}(X)=\operatorname{coker}\eta^A_{s,r,X}.
\]
These vanish simultaneously exactly when the comparison is an isomorphism.

\begin{definition}[Residual promotion]
If $X$ is a seed-generated object and $\Delta^A_{s,r}(X)$ is nonzero, the operation
\[
X\longmapsto\Delta^A_{s,r}(X)
\]
is called \emph{residual promotion} when the residual is added to the active seed set.
\end{definition}

This creates the recursive pattern
\[
\text{seed}
\to
\text{mixed square}
\to
\text{residual}
\to
\text{new seed}
\to
\text{expanded zoo}.
\]

\section{Normal forms for mixed words}

To make the two-axis calculus computational, mixed transformation-resolution expressions should admit controlled rewriting.

\subsection{Typed word system}

Consider a typed alphabet containing transformation symbols $T$ and resolution symbols $C$. A mixed word is a composable expression such as
\[
W=T_1C_1T_2T_3C_2T_4.
\]
Suppose exchange rules have the form
\[
TC\;\Rightarrow\;C\,T',
\]
where $T'$ is the correctly reindexed transformation at the new resolution.

\begin{definition}[Inversion number]
For a word $W$, let $\iota(W)$ be the number of ordered pairs $(T,C)$ in which a transformation symbol occurs to the left of a resolution symbol when the chosen normal form requires all resolution symbols first.
\end{definition}

\begin{lemma}[Exchange termination]\label{lem:termination}
Every exchange rewrite
\[
TC\Rightarrow CT'
\]
reduces $\iota$ by one. Therefore every sequence of such rewrites terminates.
\end{lemma}

\begin{proof}
The rewritten adjacent pair contributes one inversion before the rewrite and none afterward. The relative order of every other transformation-resolution pair is unchanged.
\end{proof}

\begin{theorem}[Exact two-axis normal form]\label{thm:normal-form}
Assume the exchange rewriting system is locally confluent. Then every well-typed mixed word has a unique normal form up to the equality notion of the rewriting system. The normal form has all resolution operations grouped before all within-resolution transformations.
\end{theorem}

\begin{proof}
By \cref{lem:termination}, the system terminates. A terminating locally confluent rewriting system is confluent by Newman's lemma. Therefore every word reduces to a unique normal form.
\end{proof}

\status{proved here from standard rewriting theory under local confluence}

The nontrivial application problem is to prove local confluence for a concrete Komorebi exchange signature.

\subsection{Split residual expansion}

Exact exchange is too restrictive for many applications. A first rigorous residual model can be obtained in an additive setting.

\begin{definition}[Split residual exchange]
An exchange step is \emph{split residual} if there is a natural decomposition
\[
TC(X)\cong CT'(X)\oplus\Delta_{T,C}(X).
\]
\end{definition}

\begin{theorem}[Finite split residual expansion]\label{thm:split-expansion}
Assume every exchange step in a finite mixed word admits a split residual exchange and all operators distribute over finite direct sums. Then the value of the word on a seed $s$ admits a finite direct-sum decomposition
\[
W(s)
\cong
W_{\mathrm{nf}}(s)
\oplus
\bigoplus_{\gamma\in\Gamma(W)}R_\gamma(s),
\]
where $W_{\mathrm{nf}}$ is the exact normal-form branch and each $R_\gamma$ contains at least one elementary residual $\Delta_{T,C}$.
\end{theorem}

\begin{proof}
Induct on the inversion number $\iota(W)$. If $\iota(W)=0$, take $\Gamma(W)=\varnothing$. Otherwise choose an adjacent inversion $TC$. Replace it by
\[
CT'\oplus\Delta_{T,C}.
\]
By distributivity, $W(s)$ splits into two terms. The principal term has inversion number one smaller. Apply the induction hypothesis to it. The residual branch already contains an elementary residual and can itself be normalized recursively. Finiteness follows because every principal exchange reduces inversion number and the original word is finite.
\end{proof}

\begin{definition}[Normalization path]
Let $\mathcal S$ be a stable $\infty$-category, or a stable enhancement with functorial cofibers. A \emph{normalization path} for a mixed word $W$ evaluated on $s$ is a finite chain
\[
X_0=W(s)\xrightarrow{f_1}X_1\xrightarrow{f_2}\cdots\xrightarrow{f_m}X_m=W_{\mathrm{nf}}(s),
\]
where every $f_i$ is obtained by applying an exact surrounding context to one elementary exchange morphism. Write
\[
D_i:=\Cofib(f_i),
\qquad
D_{\mathrm{tot}}:=\Cofib(X_0\to X_m).
\]
\end{definition}

\begin{theorem}[Stable residual filtration along a normalization path]\label{thm:stable-filtration}
For every finite normalization path in a stable enhancement there is a canonical finite filtration
\[
0=F_0\longrightarrow F_1\longrightarrow\cdots\longrightarrow F_m=D_{\mathrm{tot}}
\]
with cofiber sequences
\[
F_{i-1}\longrightarrow F_i\longrightarrow D_i.
\]
Consequently the associated graded pieces of the total normalization defect are precisely the transported elementary exchange residuals $D_i$.
\end{theorem}

\begin{proof}
Define
\[
F_i:=\Cofib(X_0\longrightarrow X_i).
\]
Then $F_0\simeq0$ and $F_m=D_{\mathrm{tot}}$. For the composable maps
\[
X_0\longrightarrow X_{i-1}\xrightarrow{f_i}X_i,
\]
the octahedral cofiber sequence gives
\[
\Cofib(X_0\to X_{i-1})
\longrightarrow
\Cofib(X_0\to X_i)
\longrightarrow
\Cofib(f_i).
\]
Thus
\[
F_{i-1}\longrightarrow F_i\longrightarrow D_i
\]
is a cofiber sequence. Since $f_i$ is obtained by placing an elementary exchange morphism inside an exact context, exactness identifies $D_i$ with the same context applied to the elementary exchange cofiber. Iterating gives the stated filtration.
\end{proof}

\status{proved for a chosen normalization path in a stable enhancement}

\begin{corollary}[Source-to-normal-form residual control]\label{cor:source-normal-control}
There is a cofiber sequence
\[
D_{\mathrm{tot}}[-1]\longrightarrow W(s)\longrightarrow W_{\mathrm{nf}}(s),
\]
and $D_{\mathrm{tot}}[-1]$ carries the shifted filtration with successive pieces $D_i[-1]$. Hence the failure of the original word to equal its normal form is an iterated extension of shifted elementary exchange residuals.
\end{corollary}

The theorem is path-relative. If two normalization paths are related by coherent higher exchange cells, their total comparison maps may agree up to homotopy, but path-independence of the induced filtration is a separate problem. In an ordinary triangulated category without an enhancement, the octahedral axiom gives existence of such filtrations but not a canonical functorial choice.

\begin{problem}[Order-independent residual filtration]\label{prob:path-independent}
Give coherence hypotheses under which different normalization paths induce canonically equivalent filtered total residuals, not merely equivalent normal-form objects.
\end{problem}

\section{Obstruction theory}

Residual and obstruction should not be identified. A residual is a structured discrepancy. An obstruction is a discrepancy that cannot be removed by the admissible calculus.

\begin{definition}[Killable residual]
Let $\delta\in\Zoo(\Seeds_0)$. Call $\delta$ \emph{killable} if there exists an admissible finite constructor word $W$ such that
\[
W(\delta)\simeq0.
\]
\end{definition}

\begin{definition}[Persistent obstruction]
A residual seed $\delta$ is a \emph{persistent obstruction} relative to a constructor zoo if it is not killable by any admissible finite word.
\end{definition}

This definition is intentionally operational. More refined applications should replace it by an obstruction class in a quotient, derived functor, cohomology group, or localization.

A first taxonomy is:
\[
\boxed{
\Omega=(\Omega_R,\Omega_T,\Omega_{TR},\Omega_G)
}
\]
with:
\begin{itemize}
    \item $\Omega_R$: no admissible refinement exists;
    \item $\Omega_T$: a required within-resolution transformation does not exist;
    \item $\Omega_{TR}$: exchange cannot be made coherent or its residual persists;
    \item $\Omega_G$: locally compatible data fail to reconstruct globally.
\end{itemize}

\subsection{Universal obstruction quotient in the stable calculus}

Let $\mathcal Z$ be a small idempotent-complete stable category representing a stable seed zoo, and let $\mathscr M$ be a collection of exact constructor endofunctors of $\mathcal Z$ containing the identity and closed under composition.

\begin{definition}[Operationally killable subcategory]
Define
\[
\mathcal K_{\mathscr M}
:=
\{X\in\mathcal Z:\exists W\in\mathscr M,\;W(X)\simeq0\}.
\]
Let
\[
\mathcal N_{\mathscr M}
:=
\operatorname{Thick}(\mathcal K_{\mathscr M})
\]
be the smallest thick stable subcategory containing all operationally killable objects.
\end{definition}

The thick closure is necessary in general: two objects may be killed by unrelated constructor words with no single constructor word killing an extension of them.

\begin{definition}[Stable obstruction category and obstruction class]\label{def:obstruction-quotient}
The \emph{stable obstruction category} is the Verdier quotient
\[
\mathcal O_{\mathscr M}
:=
\mathcal Z/\mathcal N_{\mathscr M},
\]
with quotient functor
\[
q:\mathcal Z\longrightarrow\mathcal O_{\mathscr M}.
\]
For $X\in\mathcal Z$, define its \emph{obstruction class}
\[
\Obs_{\mathscr M}(X):=q(X).
\]
\end{definition}

\begin{theorem}[Universal stable obstruction quotient]\label{thm:obstruction-quotient}
The quotient $q$ is exact and satisfies:
\begin{enumerate}
    \item every operationally killable object maps to zero;
    \item $q(X)\simeq0$ if and only if $X\in\mathcal N_{\mathscr M}$;
    \item for every exact functor $E:\mathcal Z\to\mathcal D$ into a triangulated or stable target such that $E(K)\simeq0$ for every $K\in\mathcal K_{\mathscr M}$, there is a unique exact factorization
    \[
    \mathcal Z\xrightarrow{q}\mathcal O_{\mathscr M}\xrightarrow{\bar E}\mathcal D.
    \]
\end{enumerate}
Thus $\mathcal O_{\mathscr M}$ is universal among exact targets in which all operationally killable residuals vanish.
\end{theorem}

\begin{proof}
The Verdier quotient by the thick subcategory $\mathcal N_{\mathscr M}$ exists and the quotient functor is exact. Its kernel is exactly $\mathcal N_{\mathscr M}$. This gives (1)--(2). For (3), the kernel of an exact functor is thick. Since $E$ kills $\mathcal K_{\mathscr M}$, it kills the thick subcategory generated by that class, namely $\mathcal N_{\mathscr M}$. The universal property of the Verdier quotient gives the unique exact factorization.
\end{proof}

\status{proved in the stable/triangulated regime using the standard Verdier quotient}

\begin{example}[Why thick closure is genuinely necessary]\label{ex:killable-not-thick}
Let $\mathcal C$ be any nonzero stable category and set
\[
\mathcal Z=\mathcal C^3.
\]
Define exact idempotent endofunctors
\[
U(X_1,X_2,X_3)=(0,0,X_2\oplus X_3),
\qquad
V(X_1,X_2,X_3)=(0,0,X_1\oplus X_3).
\]
They satisfy
\[
U^2\simeq U,
\qquad
V^2\simeq V,
\qquad
U\circ V\simeq V,
\qquad
V\circ U\simeq U,
\]
so $\mathscr M=\{\id,U,V\}$ is closed under composition. Choose $0\not\simeq A\in\mathcal C$ and set
\[
X=(A,0,0),
\qquad
Y=(0,A,0).
\]
Then $U(X)\simeq0$ and $V(Y)\simeq0$, so both $X$ and $Y$ are operationally killable. However
\[
X\oplus Y=(A,A,0),
\]
and both $U$ and $V$ send this object to $(0,0,A)$, which is fixed by $U$ and $V$. Hence no constructor word kills $X\oplus Y$. Thus $\mathcal K_{\mathscr M}$ need not be additive, much less thick.
\end{example}

This example shows that replacing $\mathcal K_{\mathscr M}$ by its thick closure in \cref{def:obstruction-quotient} is forced if the quotient is required to be universal for exact functors.

\begin{definition}[Directed annihilator property]
The constructor system $\mathscr M$ has the \emph{directed annihilator property} if for every $U,V\in\mathscr M$ there exist $P,Q,W\in\mathscr M$ and natural equivalences
\[
W\simeq P\circ U\simeq Q\circ V.
\]
Equivalently, every pair of constructor words has a common right multiple inside $\mathscr M$.
\end{definition}

\begin{theorem}[When operational killability is already thick]\label{thm:killable-thick}
If $\mathscr M$ consists of exact endofunctors and has the directed annihilator property, then $\mathcal K_{\mathscr M}$ is a thick stable subcategory. Hence
\[
\mathcal N_{\mathscr M}=\mathcal K_{\mathscr M}
\]
and
\[
\boxed{
\Obs_{\mathscr M}(X)\simeq0
\iff
X\text{ is operationally killable}.}
\]
\end{theorem}

\begin{proof}
The zero object is killed by every exact constructor. If $U(X)\simeq0$, exactness gives $U(X[1])\simeq U(X)[1]\simeq0$, so the class is closed under shifts. It is closed under retracts because exact functors preserve retract diagrams.

For cofibers, suppose $X$ is killed by $U$ and $Y$ is killed by $V$. Choose a common right multiple
\[
W\simeq P\circ U\simeq Q\circ V.
\]
Then $W(X)\simeq0\simeq W(Y)$. For any morphism $f:X\to Y$, exactness gives
\[
W(\Cofib(f))\simeq\Cofib(Wf)\simeq\Cofib(0\to0)\simeq0.
\]
Thus the killable objects form a stable subcategory closed under retracts, hence a thick subcategory.
\end{proof}

\begin{example}[Idempotent projection model]
Let $\mathcal Z=\mathcal A\times\mathcal B$ for stable categories $\mathcal A,\mathcal B$, and let
\[
L(a,b)=(a,0).
\]
With constructor monoid generated by the idempotent exact functor $L$, an object is operationally killable exactly when it lies in $0\times\mathcal B$. The obstruction quotient is equivalent to $\mathcal A$, and
\[
\Obs(a,b)\simeq a.
\]
Thus the quotient literally retains the component that cannot be annihilated by the admissible constructor.
\end{example}

\begin{theorem}[Functorial transport of obstruction]\label{thm:obstruction-transport}
Let $(\mathcal Z_A,\mathscr M_A)$ and $(\mathcal Z_B,\mathscr M_B)$ be stable obstruction systems, and let
\[
\Lambda:\mathcal Z_A\longrightarrow\mathcal Z_B
\]
be exact. If
\[
\Lambda(\mathcal K_{\mathscr M_A})\subseteq\mathcal N_{\mathscr M_B},
\]
then $\Lambda$ induces an exact functor
\[
\bar\Lambda:\mathcal O_{\mathscr M_A}\longrightarrow\mathcal O_{\mathscr M_B}.
\]
If $\Lambda$ is an equivalence and carries $\mathcal N_{\mathscr M_A}$ onto $\mathcal N_{\mathscr M_B}$, then $\bar\Lambda$ is an equivalence. In that case obstruction vanishing is preserved and reflected.
\end{theorem}

\begin{proof}
The hypothesis and exactness imply
\[
\Lambda(\mathcal N_{\mathscr M_A})\subseteq\mathcal N_{\mathscr M_B},
\]
because exact functors preserve the operations used to generate a thick subcategory. Therefore $q_B\Lambda$ kills $\mathcal N_{\mathscr M_A}$, so the universal property in \cref{thm:obstruction-quotient} gives $\bar\Lambda$. If $\Lambda$ is an equivalence identifying the two thick kernels, a quasi-inverse descends in the same way and yields a quasi-inverse to $\bar\Lambda$.
\end{proof}

This theorem gives a precise version of the intended local-global principle: a correspondence need not merely preserve local invariants; it can transport the quotient in which removable residual structure has been annihilated.

\begin{theorem}[Obstruction localization along normalization]\label{thm:obstruction-localization}
Let
\[
X_0\xrightarrow{f_1}X_1\xrightarrow{}\cdots\xrightarrow{f_m}X_m
\]
be a normalization path as in \cref{thm:stable-filtration}, with elementary transported residuals $D_i=\Cofib(f_i)$ and total residual
\[
D_{\mathrm{tot}}=\Cofib(X_0\to X_m).
\]
Let
\[
q:\mathcal Z\to\mathcal O_{\mathscr M}
\]
be the stable obstruction quotient. Then:
\begin{enumerate}
    \item $q(D_{\mathrm{tot}})$ inherits a finite filtration whose successive cofibers are $q(D_i)$;
    \item if $q(D_i)\simeq0$ for every $i$, then
    \[
    q(D_{\mathrm{tot}})\simeq0;
    \]
    \item equivalently,
    \[
    q(D_{\mathrm{tot}})\not\simeq0
    \Longrightarrow
    \exists i,\;q(D_i)\not\simeq0.
    \]
\end{enumerate}
If the directed annihilator property of \cref{thm:killable-thick} holds, the phrase ``$q(D)\simeq0$'' can be replaced throughout by ``$D$ is operationally killable.''
\end{theorem}

\begin{proof}
The residual filtration in \cref{thm:stable-filtration} consists of cofiber sequences
\[
F_{i-1}\longrightarrow F_i\longrightarrow D_i,
\qquad
F_0\simeq0,
\qquad
F_m\simeq D_{\mathrm{tot}}.
\]
Because the Verdier quotient functor $q$ is exact, applying $q$ gives cofiber sequences
\[
q(F_{i-1})\longrightarrow q(F_i)\longrightarrow q(D_i).
\]
This proves (1). If every $q(D_i)$ vanishes, then $q(F_{i-1})\to q(F_i)$ is an equivalence for every $i$. Starting from $q(F_0)\simeq0$ yields $q(F_m)\simeq0$, proving (2). Statement (3) is its contrapositive. Under the directed annihilator hypothesis, \cref{thm:killable-thick} identifies quotient-vanishing with operational killability.
\end{proof}

\status{proved from the stable residual filtration and exactness of the obstruction quotient}

This result gives the calculus a precise diagnostic property: a persistent obstruction in a composite transformation-resolution word cannot appear from nowhere. Relative to a chosen normalization path, it must be visible in at least one elementary exchange residual after passage to the obstruction quotient.

\begin{problem}[Nonstable intrinsic obstruction]
Construct an analogue of \cref{def:obstruction-quotient} before stable/additive completion, or prove that the stable quotient is the correct universal recipient for the applications of interest.
\end{problem}

\section{Local-to-global reconstruction}

\subsection{A Cech torsor residual as a nonstable obstruction test bed}

The stable Verdier quotient is not the only setting in which the residual/repair/obstruction pattern appears. A classical torsor gives a minimal nonstable model.

\begin{theorem}[Torsor residual obstruction]\label{thm:torsor-obstruction}
Let $\mathcal A$ be a sheaf of abelian groups on a space or site, let $P$ be an $\mathcal A$-torsor, and let $\{U_i\to U\}$ be a cover on which $P$ admits local sections $s_i\in P(U_i)$. On overlaps there are unique elements
\[
a_{ij}\in\mathcal A(U_{ij})
\]
such that
\[
s_j=s_i+a_{ij}.
\]
Then:
\begin{enumerate}
    \item $a=(a_{ij})$ is a Cech $1$-cocycle;
    \item changing the local sections by a $0$-cochain $b=(b_i)$,
    \[
    s_i'=s_i+b_i,
    \]
    changes the residual by the Cech coboundary
    \[
    a'_{ij}=a_{ij}+b_j-b_i;
    \]
    \item the class
    \[
    \Omega(P;\{U_i\})=[a]\in\check H^1(\{U_i\},\mathcal A)
    \]
    is independent of the chosen local sections;
    \item $\Omega(P;\{U_i\})=0$ if and only if the local sections can be modified by admissible local repairs so that they glue to a global section of $P$ over $U$.
\end{enumerate}
\end{theorem}

\begin{proof}
On a triple overlap,
\[
s_k=s_j+a_{jk}=s_i+a_{ij}+a_{jk},
\]
while also $s_k=s_i+a_{ik}$. Freeness of the torsor action gives
\[
a_{ij}+a_{jk}=a_{ik},
\]
which is the Cech cocycle identity.

If $s_i'=s_i+b_i$, then
\[
s_j'=s_j+b_j=s_i+a_{ij}+b_j
=s_i'+(a_{ij}+b_j-b_i),
\]
so the residual changes by the coboundary $\delta b$. Hence its cohomology class is independent of the chosen local sections.

If $[a]=0$, choose $b$ so that $a+\delta b=0$. The repaired sections $s_i'=s_i+b_i$ then agree on every overlap, hence glue by the sheaf property to a global section of $P$. Conversely, if a global section exists, its restrictions have zero transition residual. Comparing those restrictions with the original $s_i$ shows that $a$ differs from zero by a coboundary, so $[a]=0$.
\end{proof}

\status{proved directly from torsor descent and Cech cohomology}

In Komorebi language,
\[
\boxed{
\text{overlap residual}
\;/\;
\text{local repair}
\;=\;
\text{global obstruction class}.}
\]
Here the residual is $a_{ij}$, admissible repair is the change $s_i\mapsto s_i+b_i$, killable residuals are coboundaries, and the obstruction quotient is the cohomology class $[a]$. This example is important because it realizes the proposed obstruction pattern before stabilization and shows that the abstract vocabulary recovers an established local-to-global obstruction mechanism rather than manufacturing one artificially.
It is specifically an obstruction to reconstructing a global \emph{section} (equivalently, to trivializing the torsor), not an obstruction to the existence of the torsor itself.

\subsection{Descent data}

Let $(\mathcal C,J)$ be a site and let $\mathscr K$ assign a Komorebi category to each object of the site, with restriction functors on overlaps.

Given a cover $\{U_i\to U\}$, local objects $X_i\in\mathscr K(U_i)$ carry descent isomorphisms
\[
\phi_{ij}:X_i|_{U_{ij}}\xrightarrow{\sim}X_j|_{U_{ij}}
\]
satisfying the cocycle condition on triple overlaps.

\begin{definition}[Effective reconstruction]
The Komorebi system has \emph{effective reconstruction} on a cover if every valid descent datum arises from a global object $X\in\mathscr K(U)$, unique up to the appropriate equivalence.
\end{definition}

This is ordinary stack-style descent expressed in the present language. The Komorebi-specific question is how transformation residuals and resolution residuals contribute to the failure of descent.

\begin{theorem}[Descent transfer lemma]\label{thm:descent-transfer}
Let $A_i$ and $B_i$ be categories attached to a finite cover, with restriction functors on overlaps. Suppose equivalences
\[
\Lambda_i:A_i\simeq B_i
\]
commute with every restriction functor up to coherent natural isomorphism. Then $\Lambda$ induces an equivalence between the categories of descent data for $A$ and $B$.
\end{theorem}

\begin{proof}
Apply $\Lambda_i$ to every local object and the coherent comparison isomorphisms to every overlap map. Because the $\Lambda_i$ are equivalences and the comparison maps satisfy coherence, cocycle conditions are preserved and reflected. Quasi-inverses to the $\Lambda_i$ define a quasi-inverse functor on descent data.
\end{proof}

\begin{corollary}[Global reconstruction transfer]\label{cor:global-transfer}
If both systems have effective descent on the cover, the equivalence of descent categories induces an equivalence between their global categories over $U$.
\end{corollary}

This is the abstract form of the principle:
\[
\boxed{
\text{local equivalence}
+
\text{restriction compatibility}
+
\text{effective descent}
\Longrightarrow
\text{global equivalence}.
}
\]

\section{Seeded local-global reciprocity blueprint}

The previous section suggests a precise architecture for Langlands-facing work without claiming that the arithmetic hypotheses have been discharged.

Let
\[
\mathbb K_{\mathrm{arith}}
\qquad\text{and}\qquad
\mathbb K_{\mathrm{aut}}
\]
be arithmetic and automorphic seeded two-axis calculi.

A strong correspondence should include:
\begin{align*}
\Lambda(TX)&\simeq T'\Lambda(X),\\
\Lambda(CX)&\simeq C'\Lambda(X),\\
\Lambda(\Delta_{T,C}X)&\simeq\Delta_{T',C'}\Lambda(X),\\
\Lambda(\Obs X)&\simeq\Obs'\Lambda(X).
\end{align*}

The last two requirements are stronger than matching ordinary local invariants. They demand that the correspondence preserve the structure of failed exchange and failed reconstruction.

\begin{conjecture}[Seeded reciprocity principle]\label{conj:seeded-reciprocity}
Suppose arithmetic and automorphic Komorebi systems satisfy:
\begin{enumerate}
    \item local equivalence at every required place and resolution;
    \item compatibility with resolution restriction;
    \item preservation of the relevant local invariants;
    \item transport of mixed residual and obstruction data;
    \item effective descent on both sides;
    \item a global separation principle ensuring uniqueness from sufficiently rich local data.
\end{enumerate}
Then the induced local equivalence extends to an equivalence of the corresponding admissible global categories.
\end{conjecture}

Under categorical formulations strong enough to make items (1)--(5) an equivalence of descent systems, the conclusion follows from \cref{cor:global-transfer}. The research content lies in proving that the arithmetic and automorphic systems really instantiate those hypotheses.

\paragraph{Interpretation.}
This reframes a possible route to local-global reciprocity as:
\[
\boxed{
\text{identify the globalization obstruction}
\quad\text{and prove that local correspondence transports it.}
}
\]
If
\[
\Lambda_*(\Omega_{\mathrm{arith}})=\Omega_{\mathrm{aut}},
\]
then vanishing of one obstruction is equivalent to vanishing of the other. Effective reconstruction then turns local correspondence into global correspondence.

\section{Toy models}

\subsection{$p$-adic congruence resolution}

Let the resolution level be $n\ge1$ and take
\[
X_n=\mathbb Z/p^n\mathbb Z.
\]
For $m\le n$, coarsening is reduction
\[
C_{n\to m}:\mathbb Z/p^n\mathbb Z\to\mathbb Z/p^m\mathbb Z.
\]
For any polynomial $f\in\mathbb Z[x]$, the induced transformation satisfies
\[
C_{n\to m}(f(x))
=
f(C_{n\to m}(x)).
\]
Thus the transformation-resolution square commutes exactly and
\[
\Delta^f_{n,m}(x)=0.
\]
This is a simple instance where arithmetic transformation is resolution-natural.

By contrast, operations involving division by $p$ need not descend uniformly. Their domain of definition and lift multiplicity depend on valuation. This gives a natural source of nontrivial resolution obstruction.

\subsection{Dyadic averaging}

Let a fine state be
\[
x=(x_1,x_2)\in\mathbb R^2
\]
and let coarse resolution record only the average
\[
C(x)=\frac{x_1+x_2}{2}.
\]
Let the fine transformation be
\[
A_1(x_1,x_2)=(ax_1,bx_2)
\]
and the coarse transformation be multiplication by $c$:
\[
A_0(y)=cy.
\]
Then
\[
C(A_1x)-A_0(Cx)
=
\frac{(a-c)x_1+(b-c)x_2}{2}.
\]
Choosing
\[
c=\frac{a+b}{2}
\]
gives
\[
C(A_1x)-A_0(Cx)
=
\frac{a-b}{4}(x_1-x_2).
\]
The defect depends only on the antisymmetric fine mode $x_1-x_2$, precisely the information erased by averaging. This is a concrete model of a mixed residual exposing structure invisible at coarse resolution.

\subsection{Exact symmetry under coarsening}

For dyadic vectors, pairwise averaging commutes with reversal up to reversal of the coarse vector. This gives an exact exchange law and hence a zero mixed residual. The contrast with the previous example illustrates that residuals depend jointly on the transformation and the resolution operator.

\section{Derived invariants for the zoo}

The two-axis formalization suggests new quantities that can be defined without changing the underlying category theory.

\subsection{Commutation depth}

For a transformation $A$ and an object $X$, define the finest resolution down to which $A$ commutes with coarsening:
\[
\kappa_A(X)
:=
\sup\left\{r:
\Delta^A_{s,r}(X)\simeq0
\text{ for the relevant }s\ge r
\right\}.
\]
A low commutation depth means the transformation is strongly resolution-sensitive.

\subsection{Residual complexity}

When a numerical size function $\|\cdot\|$ is available on the fiber, define
\[
\mathcal E_A(s,r;X):=\|\Delta^A_{s,r}(X)\|.
\]
Possible choices include dimension, rank, norm, entropy, categorical amplitude, support size, or proof-complexity measures.

\subsection{Resolution branching profile}

For a chain
\[
r_0<r_1<\cdots<r_n,
\]
record
\[
B_X=(b_{r_0\to r_1}(X),\ldots,b_{r_{n-1}\to r_n}(X)).
\]
This profile measures how hidden multiplicity emerges under refinement.

\subsection{Zoo growth spectrum}

Instead of a single growth function, separate growth by generator type:
\[
g_T(n),\qquad g_R(n),\qquad g_\Delta(n).
\]
This may distinguish theories in which most complexity comes from representation changes, resolution branching, or residual generation.

\section{Candidate axioms for a foundational version}

The present note derives its first results from ordinary category and order theory. A stronger foundational program would isolate Komorebi-specific axioms.

\begin{axiom}[Two-axis presentation]
Every admissible structured state is represented in a system
\[
\pi:\K\to\Rcat
\]
whose fibers encode transformations and whose base encodes resolution.
\end{axiom}

\begin{axiom}[Resolution compositionality]
Coarsening is functorial:
\[
C_{t\to r}=C_{s\to r}C_{t\to s}.
\]
\end{axiom}

\begin{axiom}[Exchange data]
Every designated transformation family carries a coherent comparison
\[
A_rC_{s\to r}\Longrightarrow C_{s\to r}A_s.
\]
\end{axiom}

\begin{axiom}[Residual closure]
Every noninvertible exchange comparison generates an admissible residual object that may enter the seed zoo.
\end{axiom}

\begin{axiom}[Effective reconstruction, domain-specific]
For a specified class of local systems, vanishing global obstruction is equivalent to effective global reconstruction.
\end{axiom}

The last axiom should not be imposed universally without justification. In serious applications it should be a theorem of a sheaf, stack, patching, or descent theory.

\section{Machine-native semantics}

The seed calculus can be interpreted as a search grammar. A machine state may contain
\[
\mathsf S=
(\Seeds,\mathscr T,\mathscr R,\Delta,\Omega,\mathsf{Cert}).
\]
For a target problem $P$, search can occur jointly over transformation and resolution:
\[
(T^\ast,R^\ast)
\in
\operatorname*{arg\,min}_{T,R}
\mathcal C\bigl(\Omega(T,R;P)\bigr),
\]
where $\mathcal C$ is a problem-dependent residual or obstruction cost.

This changes the computational question from
\[
\text{``Which operation solves the problem?''}
\]
to
\[
\boxed{
\text{``At which resolution and in which representation does the obstruction become simplest?''}
}
\]

A generic discovery loop is:
\[
\boxed{
\text{seed}
\to
\text{generate}
\to
\text{normalize}
\to
\text{measure residual}
\to
\text{promote residual}
\to
\text{expand zoo}
\to
\text{certify}.
}
\]

Formal proof assistants can serve as the certificate layer without requiring the generative search process itself to be trusted.


\section{Universal constructions: syntax and stable residual completion}\label{sec:universal}

This section addresses the first foundational novelty test. The result is deliberately split into a syntactic theorem, which is standard double-category machinery, and a residual-completion theorem, which is a formal consequence of the free stable envelope of an additive category. The combination gives a rigorous universal construction for the additive/stable version of the seeded two-axis calculus. It does \emph{not} by itself establish that the construction is absent from the literature under another name.

\subsection{Free seeded two-axis double syntax}

\begin{definition}[Seeded two-axis double computad]\label{def:twoaxis-computad}
A \emph{seeded two-axis double computad} $\mathbb S$ consists of a small double computad together with:
\begin{enumerate}
    \item a distinguished set $\Seeds\subseteq \mathbb S_0$ of seed objects;
    \item a designated class of horizontal generating $1$-cells, interpreted as transformation generators;
    \item a designated class of vertical generating $1$-cells, interpreted as resolution generators;
    \item designated generating $2$-cells, interpreted as exchange or compatibility cells.
\end{enumerate}
No composition is assumed at the level of the signature. A realization of $\mathbb S$ in a strict double category $\mathbb D$ is a map of double computads preserving the designated labels.
\end{definition}

The label data are algebraically passive: they distinguish generators but impose no new equations. Hence the ordinary free-double-category construction applies unchanged.

\begin{theorem}[Free seeded two-axis double category]\label{thm:free-double}
For every small seeded two-axis double computad $\mathbb S$, there exists a strict double category $\mathbb D_{\mathbb S}$ and a canonical map of seeded double computads
\[
\iota_{\mathbb S}:\mathbb S\longrightarrow U\mathbb D_{\mathbb S}
\]
such that for every strict double category $\mathbb D$, restriction along $\iota_{\mathbb S}$ induces a natural bijection
\[
\operatorname{DblCat}(\mathbb D_{\mathbb S},\mathbb D)
\cong
\operatorname{DblCptd}(\mathbb S,U\mathbb D),
\]
with the right-hand side restricted to label-preserving maps.
\end{theorem}

\begin{proof}
The category of small strict double categories is monadic over the appropriate category of double-computadic data; equivalently, a double computad generates a free strict double category. Apply the free functor to the underlying double computad of $\mathbb S$. Because the seed, transformation, resolution, and exchange designations are merely distinguished subsets of the generating sorts, the unit of the free/forgetful adjunction preserves them tautologically. The ordinary adjunction therefore restricts to the displayed label-preserving universal property.
\end{proof}

\status{proved by translation to standard free double-category machinery}

This establishes a genuine universal property for the \emph{two-axis syntax}. It is not a novelty claim: free double categories and double computads are established constructions; see Fairbanks--Shulman~\cite{FairbanksShulman2026}.

\subsection{Additive indexed systems}

The residual layer becomes canonical after passing to stable categories. Let $\mathbf{Cat}^{\mathrm{add}}$ denote the category (or $\infty$-category, according to the chosen model) of small additive categories and additive functors, and let $\mathbf{Cat}^{\mathrm{st}}$ denote small stable categories and exact functors.

\begin{definition}[Additive two-axis system]
An \emph{additive two-axis system} over a small resolution category $\Rcat$ is a functor
\[
F:\Rcat^{op}\longrightarrow \mathbf{Cat}^{\mathrm{add}}.
\]
Thus every fiber $F(r)$ is additive and every coarsening functor $C_{s\to r}$ is additive.
\end{definition}

For a small additive category $\mathcal A$, write $\operatorname{St}^{\mathrm{add}}(\mathcal A)$ for its free stable envelope. We use the established adjunction
\[
\operatorname{St}^{\mathrm{add}}:
\mathbf{Cat}^{\mathrm{add}}
\rightleftarrows
\mathbf{Cat}^{\mathrm{st}}:U,
\]
characterized by
\[
\operatorname{Fun}^{\mathrm{ex}}
\bigl(\operatorname{St}^{\mathrm{add}}(\mathcal A),\mathcal C\bigr)
\simeq
\operatorname{Fun}^{\oplus}(\mathcal A,U\mathcal C)
\]
for every stable $\mathcal C$; see Saunier--Winges~\cite{SaunierWinges2026}.

\begin{theorem}[Pointwise stable completion]\label{thm:pointwise-stable}
Let $\Rcat$ be small. Postcomposition with $\operatorname{St}^{\mathrm{add}}$ defines a functor
\[
\operatorname{St}^{\mathrm{add}}_{\Rcat}:
\operatorname{Fun}(\Rcat^{op},\mathbf{Cat}^{\mathrm{add}})
\longrightarrow
\operatorname{Fun}(\Rcat^{op},\mathbf{Cat}^{\mathrm{st}}),
\qquad
F\longmapsto \operatorname{St}^{\mathrm{add}}\circ F.
\]
It is left adjoint to pointwise forgetting
\[
U_{\Rcat}:
\operatorname{Fun}(\Rcat^{op},\mathbf{Cat}^{\mathrm{st}})
\longrightarrow
\operatorname{Fun}(\Rcat^{op},\mathbf{Cat}^{\mathrm{add}}).
\]
Equivalently, for every additive two-axis system $F$ and stable two-axis system $G$ there is a natural equivalence
\[
\operatorname{Map}
\bigl(\operatorname{St}^{\mathrm{add}}_{\Rcat}F,G\bigr)
\simeq
\operatorname{Map}
\bigl(F,U_{\Rcat}G\bigr).
\]
\end{theorem}

\begin{proof}
Any adjunction $L\dashv U$ induces an adjunction on functor categories by postcomposition. Concretely, a natural transformation
\[
\alpha:L\circ F\Rightarrow G
\]
has pointwise adjoint mates
\[
\alpha_r^{\flat}:F(r)\Rightarrow U(G(r)).
\]
Naturality of the unit and counit implies that the family $\{\alpha_r^{\flat}\}_r$ is natural in $r$. The mate construction is invertible and respects higher coherence. Applying this to $L=\operatorname{St}^{\mathrm{add}}$ gives the result.
\end{proof}

\status{proved}

This is the precise sense in which the stable version of the two-axis system has a \emph{free residual-complete envelope}: no extra choices are required to create finite cofibers in the fibers, and every additive map into a stable two-axis system extends essentially uniquely to an exact one.

\subsection{Exchange data and canonical residuals}

\begin{definition}[Coherent additive transformation family]
Let $F:\Rcat^{op}\to\mathbf{Cat}^{\mathrm{add}}$. A coherent additive transformation family $A$ consists of additive endofunctors
\[
A_r:F(r)\to F(r)
\]
and, for each $r\le s$, a natural transformation
\[
\eta^A_{s,r}:A_rC_{s\to r}\Longrightarrow C_{s\to r}A_s,
\]
satisfying the identity and transitivity coherences for chains $r\le s\le t$.
\end{definition}

Functoriality of $\operatorname{St}^{\mathrm{add}}$ extends $A_r$, $C_{s\to r}$, and $\eta^A_{s,r}$ to exact functors and natural transformations
\[
\widehat A_r,
\qquad
\widehat C_{s\to r},
\qquad
\widehat\eta^A_{s,r}:
\widehat A_r\widehat C_{s\to r}
\Longrightarrow
\widehat C_{s\to r}\widehat A_s.
\]

\begin{definition}[Universal stable mixed residual]\label{def:universal-residual}
For $X\in\operatorname{St}^{\mathrm{add}}(F(s))$, define
\[
\Delta^A_{s,r}(X)
:=
\Cofib\!\left(
\widehat\eta^A_{s,r,X}
\right)
\in
\operatorname{St}^{\mathrm{add}}(F(r)).
\]
\end{definition}

\begin{theorem}[Universal residual preservation]\label{thm:residual-preservation}
Let $G:\Rcat^{op}\to\mathbf{Cat}^{\mathrm{st}}$ be a stable two-axis system and let
\[
\Phi:\operatorname{St}^{\mathrm{add}}_{\Rcat}F\longrightarrow G
\]
be an exact natural transformation which intertwines the chosen transformation families and exchange cells. Then for every $r\le s$ and every $X$ there is a canonical equivalence
\[
\Phi_r\bigl(\Delta^A_{s,r}(X)\bigr)
\simeq
\Cofib\!\left(
\Phi_r(\widehat\eta^A_{s,r,X})
\right).
\]
In particular, if the exchange data in $G$ are denoted by $\eta^B$, then
\[
\Phi_r\bigl(\Delta^A_{s,r}(X)\bigr)
\simeq
\Delta^B_{s,r}(\Phi_sX).
\]
\end{theorem}

\begin{proof}
Each $\Phi_r$ is exact, hence preserves finite colimits and therefore cofiber sequences. Apply $\Phi_r$ to the defining cofiber sequence
\[
\widehat A_r\widehat C_{s\to r}X
\xrightarrow{\widehat\eta^A_{s,r,X}}
\widehat C_{s\to r}\widehat A_sX
\longrightarrow
\Delta^A_{s,r}(X).
\]
The result is a cofiber sequence in $G(r)$. Compatibility of $\Phi$ with the transformation and exchange data identifies its first morphism with the target exchange morphism, giving the stated equivalence.
\end{proof}

\status{proved}

Thus residual formation is not merely available in the stable completion: it is \emph{universal under exact realization}.

\subsection{Residual promotion is forced by stable closure}

\begin{definition}[Stable seeded zoo]
Fix seed sets $\Seeds_r\subseteq\Ob(F(r))$. A \emph{stable seeded sub-zoo} of $\operatorname{St}^{\mathrm{add}}_{\Rcat}F$ is a family of full stable subcategories
\[
\mathcal Z_r\subseteq\operatorname{St}^{\mathrm{add}}(F(r))
\]
such that:
\begin{enumerate}
    \item $\mathcal Z_r$ contains the images of all seeds in $\Seeds_r$;
    \item $\widehat A_r(\mathcal Z_r)\subseteq\mathcal Z_r$ for every designated transformation generator;
    \item $\widehat C_{s\to r}(\mathcal Z_s)\subseteq\mathcal Z_r$ for every $r\le s$.
\end{enumerate}
The \emph{minimal stable zoo} is the intersection of all such families.
\end{definition}

\begin{proposition}[Existence of the minimal stable zoo]\label{prop:min-stable-zoo}
The minimal stable zoo exists.
\end{proposition}

\begin{proof}
The family of entire fibers $\operatorname{St}^{\mathrm{add}}(F(r))$ is a stable seeded sub-zoo, so the defining class is nonempty. Arbitrary intersections of full stable subcategories are again full stable subcategories, and invariance under the designated exact functors is preserved under intersection. Hence the intersection is the least stable seeded sub-zoo.
\end{proof}

\begin{theorem}[Automatic residual promotion]\label{thm:auto-promotion}
Let $\mathcal Z$ be the minimal stable zoo. If $X\in\mathcal Z_s$, then every mixed residual
\[
\Delta^A_{s,r}(X)
\]
belongs to $\mathcal Z_r$.
\end{theorem}

\begin{proof}
By invariance,
\[
\widehat A_r\widehat C_{s\to r}X
\quad\text{and}\quad
\widehat C_{s\to r}\widehat A_sX
\]
are objects of $\mathcal Z_r$. Since $\mathcal Z_r$ is stable, it is closed under cofibers. Therefore the cofiber of the exchange morphism between these two objects, namely $\Delta^A_{s,r}(X)$, also lies in $\mathcal Z_r$.
\end{proof}

\status{proved}

\begin{corollary}[Redundancy of residual promotion in the stable calculus]\label{cor:promotion-redundant}
In the additive/stable two-axis calculus, ``promote every residual to a seed'' need not be imposed as an independent closure axiom. Once the zoo is stable and closed under transformation and coarsening, mixed residuals are already forced to lie in it.
\end{corollary}

This is a substantive simplification of the initial axiomatization. Residual promotion remains meaningful operationally for search and provenance, but mathematically it is redundant in the stable closure.

\subsection{What has and has not been proved}

Combining \cref{thm:free-double,thm:pointwise-stable,thm:residual-preservation,thm:auto-promotion} gives the following restricted universal construction.

\begin{theorem}[Additive/stable seeded two-axis completion]\label{thm:additive-stable-completion}
For a small additive indexed two-axis system with coherent additive transformation/exchange data and designated seeds, there is a canonical stable completion with the following properties:
\begin{enumerate}
    \item the resolution diagram extends universally from additive fibers to stable fibers;
    \item all coarsening and transformation functors become exact;
    \item every exchange morphism has a canonical mixed residual given by a cofiber;
    \item exact realizations preserve these residuals;
    \item the least stable seed zoo exists and is automatically closed under mixed residual formation.
\end{enumerate}
The stable completion is initial among stable resolution systems receiving the original additive resolution diagram; with compatible exchange data, residual preservation is automatic for exact realizations.
\end{theorem}

\begin{proof}
Items (1)--(2) are \cref{thm:pointwise-stable}; item (3) is \cref{def:universal-residual}; item (4) is \cref{thm:residual-preservation}; item (5) is \cref{prop:min-stable-zoo,thm:auto-promotion}. Initiality is the adjunction of \cref{thm:pointwise-stable}.
\end{proof}

\status{proved under the explicit additive/stable hypotheses}

The stronger nonadditive claim remains open. In particular, this theorem does not construct a free \emph{arbitrary} residual-complete double category in which residuals satisfy a universal cofiber property without assuming an additive/stable target. It does, however, close the universal-construction problem for the regime in which the present mixed residual is actually defined as a cofiber.


\section{Research program and theorem ledger}

\begin{longtable}{>{\raggedright\arraybackslash}p{0.09\textwidth} >{\raggedright\arraybackslash}p{0.47\textwidth} >{\raggedright\arraybackslash}p{0.32\textwidth}}
\toprule
ID & Target & Current status\\
\midrule
\endhead
K1 & Resolution monotonicity and nested equivalence relations & Proved in \cref{thm:resolution-monotonicity}\\
K2 & Grothendieck total category & Proved in \cref{thm:total-category}\\
K3 & Two-axis factorization & Proved in \cref{thm:factorization}\\
K4 & Resolution ultrametric for nested discrete resolutions & Proved in \cref{thm:ultrametric}\\
K5 & Least seed zoo / closure existence & Proved in \cref{thm:zoo-existence}\\
K6 & Finitary iterative zoo generation & Proved under continuity hypothesis in \cref{prop:finitary}\\
K7 & Lifting coherent transformation families to the total category & Proved in \cref{thm:lift-total}\\
K8 & Mixed residual vanishing criterion & Proved in stable fibers in \cref{prop:residual-vanish}\\
K9 & Exact two-axis normal form & Proved under local confluence in \cref{thm:normal-form}\\
K10 & Split residual expansion & Proved under split additive exchange in \cref{thm:split-expansion}\\
K11 & Stable residual filtration without splitness & Proved relative to a chosen normalization path in \cref{thm:stable-filtration}; path-independence open\\
K12 & Intrinsic obstruction object / obstruction quotient & Proved in the stable regime in \cref{thm:obstruction-quotient}; exact operational detection under \cref{thm:killable-thick}\\
K13 & Free seeded two-axis double syntax with universal property & Proved via double computads in \cref{thm:free-double}\\
K14 & Universal residual-complete two-axis construction & Proved for additive/stable indexed systems in \cref{thm:additive-stable-completion}; nonadditive form open\\
K15 & Effective reconstruction theorem for a nontrivial arithmetic instance & Open domain theorem\\
K16 & Descent transfer across a correspondence & Abstract lemma proved in \cref{thm:descent-transfer}\\
K17 & Seeded local-global reciprocity instantiation & Program theorem, arithmetic hypotheses open\\
K18 & Lean formalization of K1--K10 & Not yet implemented in this note\\
K19 & Pointwise stable completion universal property & Proved in \cref{thm:pointwise-stable}\\
K20 & Automatic residual promotion in stable zoos & Proved in \cref{thm:auto-promotion}\\
K21 & Functorial transport of obstruction quotients & Proved in \cref{thm:obstruction-transport}\\
K22 & Nonstable residual/repair obstruction test bed & Proved for abelian torsors in \cref{thm:torsor-obstruction}\\
K23 & Localization of persistent total obstruction to an elementary exchange & Proved in \cref{thm:obstruction-localization}\\
\bottomrule
\end{longtable}

\section{Immediate next formal steps}

The next version should prioritize the following in order.

\begin{enumerate}
    \item \textbf{Lean core.} Formalize a finite or small resolution poset, indexed categories in a tractable surrogate, nested equivalence, zoo closure on finite signatures, and the normal-form rewriting theorem.
    \item \textbf{Residual path independence.} The iterated-cofiber filtration now exists along each chosen normalization path. Identify higher coherence assumptions that make the filtered residual independent of normalization order.
    \item \textbf{Weak double-category refinement.} The strict free syntax is now settled. Determine whether the intended semantic applications require a pseudo/doubly weak double category, equipment, or framed bicategory, and state the corresponding coherence data.
    \item \textbf{Beyond stable residuals.} The additive/stable universal completion is now settled. Determine whether there is a useful nonadditive residual universal property that is not merely a formal term algebra.
    \item \textbf{Nonstable obstruction.} The Verdier obstruction quotient is now constructed in stable fibers. Determine whether a useful pre-stable quotient exists and when stabilization loses obstruction information.
    \item \textbf{Arithmetic test bed.} Instantiate the calculus first on $p$-adic congruence towers, ramification filtrations, and local factor data before attempting a global Langlands statement.
    \item \textbf{Higher descent test bed.} The abelian torsor case now identifies residual modulo repair with a Cech $H^1$ obstruction. Extend this to nonabelian torsors, gerbes, or a $H^2$-valued obstruction where failure of effective gluing itself, rather than failure of global trivialization, is detected.
    \item \textbf{Langlands interface.} Only after the preceding steps, define arithmetic and automorphic seed signatures and state exactly which local equivalences and descent theorems are external inputs.
\end{enumerate}

\section{Brainstorm: stronger directions}

The formalism above suggests several deeper lines of investigation.

\subsection{Resolution completion}
A resolution tower may admit an inverse limit
\[
X_\infty\simeq\varprojlim_r X_r.
\]
One can ask whether the seed zoo commutes with completion:
\[
\Zoo(X_\infty)
\stackrel{?}{\simeq}
\varprojlim_r\Zoo(X_r).
\]
Failure should itself define a completion residual.

\subsection{Transformation completion}
Repeated transformations may generate an orbit category. Completing this orbit under limits, colimits, idempotent splitting, or derived closure may produce a canonical transformation hull.

\subsection{Bi-completion}
There are then two completion operations:
\[
\widehat{(-)}_T,
\qquad
\widehat{(-)}_R.
\]
The mixed question is whether
\[
\widehat{\widehat X_T}_R
\simeq
\widehat{\widehat X_R}_T.
\]
Its defect is a completion-level mixed residual.

\subsection{Resolution spectrum}
For a fixed object $X$, define the set of resolutions at which a property $P$ becomes visible:
\[
\operatorname{Spec}_R(P;X)
=
\{r:P(C_{s\to r}X)\text{ is detectable}\}.
\]
This could behave like a support or spectrum for distinguishability.

\subsection{Residual genealogy}
Residual promotion creates a directed graph
\[
s\to\Delta_1\to\Delta_2\to\cdots.
\]
Classify whether genealogies terminate, become periodic, stabilize, or grow indefinitely. A terminating genealogy would correspond to finite obstruction resolution. Periodic genealogy may encode a persistent but structured defect.

\subsection{Seed minimality}
Given a generated zoo $\Zoo$, ask for the smallest seed set $\Seeds$ with
\[
\Zoo(\Seeds)=\Zoo.
\]
This is a mathematical compression problem. Nonuniqueness of minimal seeds could itself be meaningful.

\subsection{Dual zoos}
If the calculus carries a duality $D$, study
\[
D\Zoo(\Seeds)
\quad\text{versus}\quad
\Zoo(D\Seeds).
\]
The residual of dualization against zoo generation measures whether the generative grammar is duality-complete.

\subsection{Proof-resolution}
A theorem itself may have resolutions. At coarse resolution one stores only a proposition and certificate. At finer resolution one exposes lemmas, dependencies, computational witnesses, and failed branches. This creates a resolution geometry on proofs and may be useful for machine reasoning and proof compression.

\subsection{World-resolution}
A foundational extension can assign resolution to models or ``worlds'' themselves. Morphisms between worlds then have two axes: change of world and change of resolution within a world. This may connect the current calculus to the earlier compatibility/residual/obstruction maps between models of a foundational theory.

\section{Provisional philosophical interpretation}

The mathematics does not require the following interpretation, but the two-axis formalism sharpens it.

A structure may be characterized simultaneously by:
\begin{itemize}
    \item what transformations it admits;
    \item what resolutions distinguish its internal states;
    \item what survives both kinds of movement;
    \item what residuals appear when the two directions fail to commute;
    \item what obstructions prevent refinement, transformation, or global reconstruction.
\end{itemize}

This suggests the strengthened ontological phrase:
\[
\boxed{
\text{to remain identifiable is to remain coherent under both transformation and resolution.}
}
\]
The mathematical content of the project is not this sentence itself. The mathematical content is the attempt to make ``coherent,'' ``transformation,'' ``resolution,'' ``residual,'' and ``reconstruction'' precise enough that such a sentence can be tested rather than merely asserted.

\section{Conclusion}

The first formalization supports a concrete picture:
\[
\boxed{
\begin{array}{c}
\text{seeds generate states},\\
\text{transformations generate representations},\\
\text{resolutions generate distinguishability},\\
\text{exchange maps relate the axes},\\
\text{residuals measure failed exchange},\\
\text{residuals become new seeds},\\
\text{obstructions classify persistent failure},\\
\text{descent reconstructs compatible local structure}.
\end{array}
}
\]
The ordinary categorical core is already rigorous. The additive/stable residual layer has a universal construction: pointwise free stabilization supplies canonical cofiber residuals, exact realizations preserve them, residual promotion is automatic inside the minimal stable zoo, and finite normalization paths admit octahedral residual filtrations. The stable obstruction layer is represented by a universal Verdier quotient, with exact detection of operational killability under a directed-annihilator hypothesis and functorial transport under exact correspondences. Combining these structures yields an obstruction-localization theorem: elementary removable residuals imply a removable total normalization defect, while a persistent total defect forces at least one persistent elementary exchange residual. The abelian torsor test bed shows independently that, before stabilization, overlap residuals modulo admissible local repairs recover the classical Cech $H^1$ obstruction to a global section. The central research problem is therefore sharper: determine whether analogous obstruction quotients control genuinely difficult arithmetic reconstruction failures, whether residual filtrations are path-independent under higher coherence, and whether the stable and nonstable obstruction languages admit a common universal formulation.


\section*{References}
\addcontentsline{toc}{section}{References}
\begin{thebibliography}{9}
\bibitem{FairbanksShulman2026}
A.~D. Fairbanks and M.~Shulman,
\emph{Doubly Weak Double Categories},
Applied Categorical Structures \textbf{34} (2026), Article 26.
The paper develops double computads and records the free/monadic strict double-category machinery used in \cref{thm:free-double}.

\bibitem{SaunierWinges2026}
V.~Saunier and C.~Winges,
\emph{On exact categories and their stable envelopes},
Mathematische Zeitschrift \textbf{312} (2026), Article 28; arXiv:2502.03408.
Section 3 recalls the free stable envelope $\operatorname{St}^{\mathrm{add}}$ and its adjunction with the forgetful functor from stable to additive categories.

\bibitem{StacksVerdier}
The Stacks Project Authors,
\emph{Quotients of triangulated categories}, Section 13.6, Tags 05RA--05RJ, accessed September 2026.
The quotient construction, kernel description, and exact universal property are used in \cref{thm:obstruction-quotient,thm:obstruction-transport}; the octahedral axiom TR4 underlies \cref{thm:stable-filtration}.
\end{thebibliography}

\appendix
\section{Compact definition sheet}

\begin{align*}
\text{Resolution system:}\quad
&F:\Rcat^{op}\to\Cat.\\[0.4em]
\text{Fiber:}\quad
&\Tcat_r=F(r).\\[0.4em]
\text{Coarsening:}\quad
&C_{s\to r}:\Tcat_s\to\Tcat_r.\\[0.4em]
\text{Total category:}\quad
&\K=\int_{\Rcat}F.\\[0.4em]
\text{Resolution equivalence:}\quad
&X\sim_rY\iff C_{t\to r}X\cong C_{t\to r}Y.\\[0.4em]
\text{Lift set:}\quad
&\Lift_{r\to s}(X)=\{Y:C_{s\to r}Y\cong X\}/\cong.\\[0.4em]
\text{Exchange map:}\quad
&\eta^A_{s,r}:A_rC_{s\to r}\Rightarrow C_{s\to r}A_s.\\[0.4em]
\text{Stable residual:}\quad
&\Delta^A_{s,r}(X)=\Cofib(\eta^A_{s,r,X}).\\[0.4em]
\text{Stable obstruction quotient:}\quad
&\mathcal O_{\mathscr M}=\mathcal Z/\operatorname{Thick}(\mathcal K_{\mathscr M}).\\[0.4em]
\text{Obstruction class:}\quad
&\Obs_{\mathscr M}(X)=q(X).\\[0.4em]
\text{Seed zoo:}\quad
&\Zoo(\Seeds_0)=\mu A.\bigl(\Seeds_0\cup\Phi(A)\bigr).\\[0.4em]
\text{Resolution ultrametric:}\quad
&d_R(X,Y)=2^{-\nu(X,Y)}.\\[0.4em]
\text{Reconstruction target:}\quad
&\Omega_G=0\iff X_{\mathrm{global}}=\Recon(\{X_i\}).
\end{align*}

\section{Version ledger}

\begin{description}[leftmargin=8em,style=nextline]
\item[v0.5.0] Combines the stable residual filtration with the Verdier obstruction quotient to prove obstruction localization along any chosen normalization path: if all elementary exchange residuals vanish in the obstruction quotient, so does the total normalization residual; contrapositively, a persistent total obstruction forces at least one persistent elementary exchange obstruction. Under directed annihilators the statement upgrades from quotient-vanishing to operational killability.
\item[v0.4.0] Adds a direct nonstable obstruction test bed: local sections of an abelian torsor produce overlap residuals, local changes of trivialization act as repairs, and the resulting Cech $H^1$ class vanishes exactly when the local sections can be repaired to reconstruct a global section. This closes the first descent test-bed obligation and shows that the residual/repair/obstruction pattern recovers a classical cohomological mechanism.
\item[v0.3.0] Proves the stable residual filtration along a chosen normalization path by iterated octahedral cofiber sequences. Constructs the universal stable obstruction quotient by the thick closure of operationally killable residuals, proves exact detection under a directed common-annihilator hypothesis, and proves functorial transport of obstruction quotients under exact correspondences. K11 and K12 are closed in these stable forms; path-independent filtration and nonstable obstruction remain open.
\item[v0.2.0] Adds the free seeded two-axis double-computad universal property and proves the pointwise free stable completion theorem for additive indexed systems. Mixed residuals become canonical cofibers, exact realizations preserve them, and residual promotion is shown to be automatic in the minimal stable zoo. K13 is closed at the strict syntactic level and K14 is closed under additive/stable hypotheses; the nonadditive residual universal construction remains open.
\item[v0.1.0] Initial continuous formalization. Introduces indexed resolution categories, total Grothendieck category, factorization, resolution ultrametric, seed zoo closure, exchange transformations, stable residuals, exact normal forms, split residual expansion, obstruction taxonomy, descent transfer, and a seeded local-global reciprocity blueprint.
\end{description}

\end{document}
